% !PRINTED 4
\weberpage{4}
%<<CONT
prime number, or a higher power of one of the prime numbers
$a,b,c,\ldots$, would have to divide $m$.

\begin{center}
\rule{0.18\linewidth}{0.4pt}
\end{center}

We now give an overview of the extensions of the concept of number
which are necessary in mathematics and which were gradually introduced.

By a \gesperrt{manifold or set}, or by the abbreviated sign
$\mathfrak M$, we understand a system of objects or elements of any
kind, which is so delimited and complete in itself that, for every
arbitrary object, it is completely determined whether it belongs to
the system or not, regardless of whether we are able, in each
particular case, actually to make the decision or not.

A set is called \gesperrt{ordered} when, of any two different elements,
one that is fully determined in itself always counts as the greater,
and in such a way that from $a>b$, $b>c$ one always has $a>c$. If
$a>b$ and $b>c$, or $a>b>c$, we say that $b$ lies between $a$ and
$c$.% [source-original sic] The print repeats the descending case; the logically expected ascending alternative is not silently substituted.

The natural numbers form an ordered set; between two successive
elements of it there lies no further element. Such a manifold is called
\gesperrt{discrete}. An ordered set with the property that between any
two elements further elements can always be found is called
\gesperrt{dense}. One can form a dense set by combining the natural
numbers in pairs and regarding these pairs as elements of a set. These
pairs shall be called \gesperrt{fractions} and denoted by $m:n$ or
$m/n$, and two such fractions $m:n$ and $m':n'$ are set equal when
$mn'=nm'$. If all fractions equal to one another are combined into one
element, one obtains a manifold which is ordered once one further
stipulates that $m:n$ is greater than $m':n'$ when $mn'>nm'$. That this
manifold is dense is seen as follows: if $\mu=m:n$, $\mu'=m':n'$ are
two fractions and $\mu>\mu'$, then, with $h$ arbitrary, one can put
\[
 \mu={hmn'\over hnn'},
 \qquad
 \mu'={hm'n\over hnn'}
\]
%CONT>>
